% ==============================================================================
% ================================ AUFGABE 1 ===================================
% ==============================================================================

%\chapter{Aufgabe 1: Boolesche Algebra}
%\label{chap:aufgaben}
%\thispagestyle{fancy}

\section{Aufgabe 1: Boolesche Algebra}
In diesem Kapitel werden die beiden gegebenen booleschen Gleichungen schrittweise vereinfacht. 
Alle Zwischenschritte werden angegeben und mit den angewandten Gesetzen der Booleschen Algebra kommentiert, so wie Fleury, M. Otten \cite{2} es beschreibt.
% ======================================= 1a =====================================
\subsection{Aufgabe 1a}
\begin{equation}
    \boldsymbol{y}_1 = \overline{ \overline{(\overline{\boldsymbol{x}_1} \cdot \overline{\boldsymbol{x}_2} \cdot \overline{\boldsymbol{x}_4})} 
    \cap \overline{(\overline{\boldsymbol{x}_1} \cup \overline{\boldsymbol{x}_2} \cup \overline{\boldsymbol{x}_3})} } \tag{\textit{Angabe}}
\end{equation}
\subsubsection{Vereinfachung}
Zunächst wird die äußere Negation mit dem Gesetz von De Morgan aufgelöst:
\begin{align}
    y_1 &= \overline{ A \cap B } \quad \text{mit} \quad 
    A = \overline{(\overline{x_1} \cdot \overline{x_2} \cdot \overline{x_4})},\; 
    B = \overline{(\overline{x_1} \cup \overline{x_2} \cup \overline{x_3})} \nonumber \\
    &= \overline{A} \cup \overline{B} \tag{\textit{De Morgan}}
\end{align}
Nun werden die doppelten Negationen aufgelöst:
\begin{align}
    \overline{A} &= \overline{ \overline{(\overline{x_1} \cdot \overline{x_2} \cdot \overline{x_4})} } 
    = \overline{x_1} \cdot \overline{x_2} \cdot \overline{x_4} \tag{\textit{Doppelte Negation}} \\
    \overline{B} &= \overline{ \overline{(\overline{x_1} \cup \overline{x_2} \cup \overline{x_3})} } 
    = \overline{x_1} \cup \overline{x_2} \cup \overline{x_3} \tag{\textit{Doppelte Negation}} \\
    y_1 &= (\overline{x_1} \cdot \overline{x_2} \cdot \overline{x_4}) 
    \cup (\overline{x_1} \cup \overline{x_2} \cup \overline{x_3}) \tag{\textit{Zwischen Ergebnis}}
\end{align}
Der Term $\overline{x_1} \cdot \overline{x_2} \cdot \overline{x_4}$ ist eine Teilmenge 
von $\overline{x_1} \cup \overline{x_2} \cup \overline{x_3}$, da er die Konjunktion 
dreier Literale ist, während die zweite Klammer die Disjunktion dieser Literale darstellt. 
Somit gilt das Absorptionsgesetz:
\begin{equation}
    y_1 = \overline{x_1} \cup \overline{x_2} \cup \overline{x_3} \tag{\textit{Absorptionsgesetz}}
\end{equation}
Die endgültige vereinfachte Gleichung lautet:
\begin{equation}
    \boxed{ \boldsymbol{y}_1 = \overline{\boldsymbol{x}_1} \cup \overline{\boldsymbol{x}_2} \cup \overline{\boldsymbol{x}_3} } \tag{\textit{Ergebnis}}
\end{equation}
% ======================================= 1b =====================================
\subsection{Aufgabe 1b}
\begin{equation}
    \boldsymbol{y}_2 = \overline{(\overline{\boldsymbol{x}_1} \cdot \boldsymbol{x}_2 \cdot \overline{\boldsymbol{x}_3} \cup \overline{(\boldsymbol{x}_1 \cup \boldsymbol{x}_2 \cup \boldsymbol{x}_3)})} \cap (\boldsymbol{x}_1 \cup \overline{\boldsymbol{x}_2}) \tag{\textit{Angabe}}
\end{equation}
\subsubsection*{Vereinfachung}
Zunächst wird die Negation der gesamten ersten Klammer mit De Morgan aufgelöst:
\begin{align}
    y_2 &= \overline{ A \cup B } \cap (x_1 \cup \overline{x_2}) \quad \text{mit} \quad
    A = \overline{x_1} \cdot x_2 \cdot \overline{x_3},\; 
    B = \overline{(x_1 \cup x_2 \cup x_3)} \nonumber \\
    &= (\overline{A} \cap \overline{B}) \cap (x_1 \cup \overline{x_2}) \tag{\textit{De Morgan}}
\end{align}
Die doppelten Negationen werden aufgelöst:
\begin{align}
    \overline{A} &= \overline{(\overline{x_1} \cdot x_2 \cdot \overline{x_3})} 
    = x_1 \cup \overline{x_2} \cup x_3 \tag{\textit{De Morgan}} \\
    \overline{B} &= \overline{ \overline{(x_1 \cup x_2 \cup x_3)} } 
    = x_1 \cup x_2 \cup x_3 \tag{\textit{Doppelte Negation}}
\end{align}
Somit wird der erste Faktor:
\begin{equation}
    \overline{A} \cap \overline{B} 
    = (x_1 \cup \overline{x_2} \cup x_3) 
    \cap (x_1 \cup x_2 \cup x_3)
    = x_1 \cup x_3 \tag{\textit{Absorption}}
\end{equation}
Eingesetzt in die ursprüngliche Gleichung ergibt sich:
\begin{equation}
    y_2 = (x_1 \cup x_3) 
    \cap (x_1 \cup \overline{x_2}) \tag{\textit{Eingesetzt}}
\end{equation}
Da $x_1$ in beiden Klammern vorkommt, kann das Absorptionsgesetz 
angewendet werden:
\begin{equation}
    y_2 = x_1 \cup (x_3 \cap \overline{x_2}) \tag{\textit{Absorptionsgesetz}}
\end{equation}
Die endgültige vereinfachte Gleichung lautet:
\begin{equation}
    \boxed{ \boldsymbol{y}_2 = \boldsymbol{x}_1 \cup (\boldsymbol{x}_3 \cap \overline{\boldsymbol{x}_2}) } \tag{\textit{Ergebnis}}
\end{equation}

% ==============================================================================
% ================================ AUFGABE 1 ===================================
% ==============================================================================